% Scientific AI: empirically verified replacement text for Figures 2–3.
% Insert into the corresponding section only after matching the original
% document labels/citations and confirming the generated figure assets exist.
\paragraph{Scientific inference from observed sunspot activity.}
Long-term annual sunspot observations offer a quantitative example of a measured natural phenomenon whose temporal variability can be investigated by data-driven methods. The original National Geophysical Data Center record distributed with \texttt{statsmodels} contains 309 annual observations from 1700 through 2008. Figure~\ref{fig:sci_ai_real_sunspots} presents the original measurements, their 11-year running average, an observed lag-1 scatter plot, and a held-out forecasting evaluation. Each plot derives from the original recorded values or from an actually fitted regression model rather than from simulated observations.

For the forecasting panel, 12 recorded annual values were used to predict the next observed annual count. The chronological partition included 178 training windows (target years 1712--1889), 60 validation windows (1890--1949), and 59 held-out testing windows (1950--2008). Training-only standardization was used. The Ridge penalty parameter was selected by minimum validation RMSE over a prespecified five-value grid, yielding $\alpha=10$. Test RMSE was 19.38 annual sunspots, versus 33.18 for the one-year persistence baseline. These values demonstrate performance for a single partition of this particular observational record, not a general forecasting guarantee.

\begin{figure}[htbp]
\centering
\includegraphics[width=0.94\linewidth]{figure_02_evolution_scientific_discovery.pdf}
\caption{Scientific inference from real annual sunspot measurements. (a) Original NGDC annual observations for 1700--2008. (b) Observed annual values and their centered 11-year rolling average. (c) Original successive-year observation pairs. (d) Actual held-out 1950--2008 predictions from validation-selected Ridge regression compared with measured sunspot counts and a persistence baseline. No simulated or invented measurements are used.}
\label{fig:sci_ai_real_sunspots}
\end{figure}

\paragraph{Scientific classification from measured biomedical features.}
The Wisconsin Diagnostic Breast Cancer dataset records 30 quantitative characteristics of digitized fine-needle aspirate cell-nucleus images in 569 cases. This constitutes authentic biomedical measurement data, although the results below describe statistical classification rather than clinical deployment. Figure~\ref{fig:sci_ai_real_wdbc} connects observed feature distributions and empirical inter-feature correlations with an actually fitted logistic regression classifier and its held-out outputs.

A stratified partition assigned 341 cases to training, 114 to validation, and 114 to held-out testing. A training-only fitted standard scaler preceded logistic regression; the inverse regularization parameter $C$ was selected by malignant-class validation F1 from values 0.01, 0.1, 1, 10, and 100. The selected model used $C=1$ and achieved a measured held-out accuracy of 98.25\%, malignant-class F1 of 97.56\%, and area under the ROC curve of 0.9957. The confusion matrix contained 72 true benign, zero benign-to-malignant errors, two malignant-to-benign errors, and 40 correctly classified malignant cases. Results are specific to this dataset and split.

\begin{figure}[htbp]
\centering
\includegraphics[width=0.94\linewidth]{figure_03_evolution_scientific_ai.pdf}
\caption{An observation-grounded scientific machine-learning workflow using authentic Wisconsin Diagnostic Breast Cancer data. (a) Recorded nuclear radius and texture for training cases. (b) Measured correlation coefficients among the first eight training-set features. (c) Actual validation F1 values used for regularization selection. (d) Classifier probabilities on 114 held-out cases. (e) ROC computed from the same held-out predictions (AUC 0.9957). The panels contain neither synthetic observations nor fabricated model scores.}
\label{fig:sci_ai_real_wdbc}
\end{figure}
